\section{Introduction and Related Work} \label{sec:intro}

    Transformers have demonstrated remarkable success across a wide range of natural language processing (NLP) tasks, such as text classification, summarization, and machine translation. Nevertheless, they still face significant challenges when asked to perform \emph{multi-step} or \emph{multi-hop} factual reasoning, particularly in real-world scenarios where knowledge is both vast and sparsely distributed. A key reason for this difficulty lies in the model’s tendency to memorize rather than \emph{generalize} -- a problem that becomes acute in knowledge-intensive tasks with insufficiently rich data distributions.

    \begin{figure}[h]
        \centering
        \includesvg[width=0.5\textwidth]{introduction_accuracy.svg}
         \vspace{-24pt} % Reduce space between SVG and caption
        \caption{Average accuracy on 2WikiMultiHopQA for comparison task. Despite GPT2-small being a model with 124 million parameters, grokked version achieves almost 100\% accuracy, beating the most recent gpt-4o and o1-mini models.}
        \label{fig:introduction_accuracy}
    \end{figure}
    
    \paragraph{From Toy Grokking to Real-World Data.} 
        Recent work on \emph{grokking} \citep{power2022grokking} has shown that, under certain conditions, overparameterized neural networks suddenly transition from pure memorization to near-perfect generalization after long training. Early studies have typically focused on highly controlled, synthetic tasks such as modular arithmetic or simplified algorithmic datasets. In these \emph{toy} settings, the number of \emph{inferred facts} (i.e., multi-step or composed patterns) can be systematically increased until a threshold ratio~$\phi$ is reached, at which point a ``generalizing circuit'' emerges in the model \citep{belkin2019reconciling, nakkiran2020deepdouble, thilak2022slingshot, humayun2023dynamics, nanda2023progress}. 
    
        However, real-world datasets present a stark contrast: factual knowledge is \emph{extremely sparse} and often scattered across incomplete or noisy knowledge graphs. Thus, the core challenge is to ensure that there are enough higher-order \emph{inferred facts} in relation to the \emph{atomic facts} (direct statements) to enable the internal circuit-formation process that grokking requires. Put differently, in real-world scenarios, one cannot trivially guarantee a sufficiently large ratio $\phi$ between multi-step (inferred) facts and single-hop (atomic) facts. Our work addresses precisely this obstacle by proposing a data-synthesis strategy that augments and re-balances real-world knowledge bases.
    
    \paragraph{Multi-hop Question Answering (QA).}
        2WikiMultiHopQA \citep{yang_hotpotqa_2018, ho_2WikiMultiHopQA_2020, trivedi_MuSiQue_2022} is particularly well-suited to assess multi-step factual reasoning. This dataset contains Wikipedia-based queries that require retrieving and combining multiple pieces of evidence (spread across different pages or paragraphs) before producing an answer. This property aligns closely with our focus on multi-hop reasoning and underscores the need for implicit reasoning capability: in many cases, a system must link and traverse \emph{several} factual nodes -- e.g., ``Michelle is the wife of Obama'' and ``Michelle was born in 1964'' -- to arrive at a final conclusion (e.g., about Michelle's birth year or other derived facts). Moreover, they reflect a large, complex knowledge graph with real-world entities, ambiguous language, and long-tail relations -- all of which make them an archetypal testbed for factual reasoning at scale.
    
    \paragraph{Grokking and Its Role in Transformer Generalization.}
        The concept of grokking, introduced in \citet{power2022grokking}, demonstrated that neural networks could learn not just superficial patterns but also deeper, generalizable reasoning mechanisms under prolonged training with suitable inductive biases. Subsequent work has linked grokking to double descent \citep{belkin2019reconciling, nakkiran2020deepdouble}, the geometry of deep-network loss landscapes \citep{davies2023unifying}, and weight decay \citep{pezeshki2022multiscale, nanda2023progress}, suggesting that the right regularization can encourage the emergence of \emph{generalizing circuits}. These circuits -- once formed -- enable out-of-distribution reasoning that surpasses naive memorization \citep{varma_circuit_efficiency_2023, liu2023omnigrok}.
    
    \paragraph{Gap in the Literature.}
        Despite extensive research on \emph{knowledge graph completion} \citep{liu_joint_2022} and multi-hop question answering via retrieval-based methods \citep{yang_hotpotqa_2018, ho_2WikiMultiHopQA_2020}, very few studies have examined whether the \emph{internal grokking phenomenon} can be harnessed for implicit multi-hop reasoning in a \emph{real-world} textual setting. Most prior approaches either:
        \begin{itemize}
            \item \textbf{Focus on toy tasks}: e.g., modular addition or synthetic math problems \citep{power2022grokking, nanda2023progress, wang_grokked_2024}.
            \item \textbf{Rely on explicit prompting or chain-of-thought}: where intermediate reasoning steps must be spelled out in external text \citep{plaat2024reasoning}, rather than learned as an implicit circuit.
            \item \textbf{Use standard graph-completion architectures}: e.g., GNN-based solutions to augment partial knowledge graphs \citep{liu_joint_2022}, which do not necessarily yield late-phase \emph{internal} circuit formation or sudden generalization.
        \end{itemize}
        To the best of our knowledge, no existing work has used a \emph{grokking-based} approach to demonstrate how a Transformer can \emph{implicitly} discover multi-hop reasoning skills on large-scale factual data. 
    
    \paragraph{Contributions.}
        In this paper, we bridge that gap through the following contributions:
    
        \begin{itemize}
            \item We incorporate a targeted \emph{data synthesis} procedure to ensure sufficiently large $\phi$ for each relation, thereby unlocking the potential for \emph{internal generalization circuits} to form in real-world Wikipedia-based tasks.
            \item We show that \emph{even factually incorrect synthetic data can boost the ratio of inferred to atomic facts}, often strengthening rather than harming logical consistency. 
            \item  Our experiments on 2WikiMultiHopQA confirm that once the ratio $\phi$ surpasses a certain threshold, grokking emerges -- enabling Transformers to perform complex multi-step reasoning \emph{without} explicit chain-of-thought prompts or elaborate external scaffolding.
        \end{itemize}
        
        In the following sections, we detail the mathematical basis of our data augmentation strategy, the empirical setups on 2WikiMultiHopQA, and the new insights gained about implicit multi-hop reasoning. Our findings show that \emph{grokking is not an artifact confined to contrived toy datasets} but a powerful mechanism that, with suitable data distribution adjustments, can be harnessed for real-world factual reasoning at scale.
    